% =====================================================================
%  תרגיל בית 2 --- אחרי הרצאה 2 (הסקה בייסיאנית, MCMC, סופרנובות)
%  קומפילציה: lualatex hw2.tex
% =====================================================================
\documentclass[11pt,a4paper]{article}
% =====================================================================
%  Shared preamble — lecture notes "Introduction to Cosmology" (9 h)
%  Compile:  lualatex cosmo_lectures.tex   (twice, for the TOC)
% =====================================================================
\usepackage[a4paper,top=2.2cm,bottom=2.2cm,left=2.2cm,right=2.2cm]{geometry}
\usepackage{fontspec}
\usepackage{amsmath,amssymb}
\usepackage[bidi=basic]{babel}
\babelprovide[import,main]{hebrew}
\babelprovide[import]{english}

% ---- fonts: Culmus (Overleaf/Linux) -> Windows David -> FreeSerif ----
\IfFontExistsTF{David CLM}{%
  \babelfont{rm}{David CLM}\babelfont{sf}{Simple CLM}\babelfont{tt}{Miriam Mono CLM}%
}{\IfFontExistsTF{David}{%
  \babelfont{rm}[ItalicFeatures={FakeSlant=0.15}]{David}\babelfont{sf}{Arial}\babelfont{tt}{Courier New}%
}{\IfFileExists{/mnt/c/Windows/Fonts/david.ttf}{%
  \babelfont{rm}[Path=/mnt/c/Windows/Fonts/,Extension=.ttf,BoldFont=davidbd,ItalicFont=david,
    ItalicFeatures={FakeSlant=0.15}]{david}%
  \babelfont{sf}[Path=/mnt/c/Windows/Fonts/,Extension=.ttf,BoldFont=arialbd]{arial}%
  \babelfont{tt}[Path=/mnt/c/Windows/Fonts/,Extension=.ttf]{cour}%
}{%
  \babelfont{rm}{FreeSerif}\babelfont{sf}{FreeSans}\babelfont{tt}{FreeMono}%
}}}
\babelfont[english]{rm}{Latin Modern Roman}
\babelfont[english]{tt}{Latin Modern Mono}

\usepackage{xcolor}
\usepackage[framemethod=default]{mdframed}
\usepackage{enumitem}
\usepackage{booktabs}
\usepackage{graphicx}
\usepackage{fancyhdr}
\usepackage{needspace}
\usepackage{listings}
\usepackage[hidelinks]{hyperref}
\lstset{basicstyle=\ttfamily\footnotesize,keywordstyle=\color{qblue}\bfseries,
  commentstyle=\color{ggreen}\itshape,frame=single,rulecolor=\color{black!30},
  backgroundcolor=\color{black!2},breaklines=true,showstringspaces=false,
  language=Python,columns=fullflexible}

\renewcommand{\labelitemi}{\ensuremath{\bullet}}
\renewcommand{\labelitemii}{\ensuremath{\circ}}
\setlist[itemize]{leftmargin=1.4em,itemsep=1pt,topsep=2pt}
\setlist[enumerate]{leftmargin=1.6em,itemsep=1pt,topsep=2pt}

% ---- colours (as in the original course) ----
\definecolor{qblue}{RGB}{20,60,120}
\definecolor{ggreen}{RGB}{20,100,70}
\definecolor{kpurple}{RGB}{95,45,110}
\definecolor{worange}{RGB}{170,90,10}

\newmdenv[linecolor=kpurple,linewidth=1pt,backgroundcolor=kpurple!4,
  innertopmargin=6pt,innerbottommargin=6pt,innerleftmargin=8pt,innerrightmargin=8pt,
  skipabove=7pt,skipbelow=7pt]{keybox}
\newmdenv[linecolor=worange,linewidth=1pt,backgroundcolor=worange!5,
  innertopmargin=6pt,innerbottommargin=6pt,innerleftmargin=8pt,innerrightmargin=8pt,
  skipabove=7pt,skipbelow=7pt]{warnbox}
\newmdenv[linecolor=ggreen,linewidth=1pt,backgroundcolor=ggreen!5,
  innertopmargin=6pt,innerbottommargin=6pt,innerleftmargin=8pt,innerrightmargin=8pt,
  skipabove=7pt,skipbelow=7pt]{goalbox}
\newmdenv[linecolor=qblue,linewidth=1.1pt,backgroundcolor=qblue!5,
  innertopmargin=6pt,innerbottommargin=6pt,innerleftmargin=8pt,innerrightmargin=8pt,
  skipabove=8pt,skipbelow=6pt]{exbox}
\newmdenv[linecolor=black!35,linewidth=0.7pt,backgroundcolor=black!3,
  innertopmargin=6pt,innerbottommargin=6pt,innerleftmargin=8pt,innerrightmargin=8pt,
  skipabove=4pt,skipbelow=8pt]{solbox}

% ---- macros ----
\newcommand{\en}[1]{\foreignlanguage{english}{#1}}
\newcommand{\Msun}{M_{\odot}}
\newcommand{\Mpc}{\mathrm{Mpc}}
\newcommand{\kpc}{\mathrm{kpc}}
\newcommand{\Gyr}{\mathrm{Gyr}}
\newcommand{\kms}{\mathrm{km\,s^{-1}}}
\newcommand{\kmsMpc}{\mathrm{km\,s^{-1}\,Mpc^{-1}}}
\newcommand{\LCDM}{\Lambda\mathrm{CDM}}

% lecture header: \lecture{number}{title}{time budget + one-line content}
\newcounter{lec}
\counterwithin*{subsection}{lec}
\renewcommand{\thesubsection}{\thelec.\arabic{subsection}}
\makeatletter\renewcommand{\fps@figure}{htbp}\makeatother
\newcommand{\lecture}[3]{%
  \clearpage\refstepcounter{lec}%
  \noindent\colorbox{qblue}{\parbox{\dimexpr\textwidth-2\fboxsep}{%
   \color{white}\bfseries\Large הרצאה #1: #2\\[3pt]\normalsize\mdseries #3}}%
  \addcontentsline{toc}{section}{הרצאה #1: #2}\par\medskip}

% exercises: numbered per lecture (1.1, 1.2, ...)
\newcounter{ex}[lec]
\newenvironment{exercise}[1]{%
  \refstepcounter{ex}\begin{exbox}\textbf{\textcolor{qblue}{תרגיל \thelec.\theex\ --- #1.}}\ }%
  {\end{exbox}}
\newcommand{\soln}[1]{\par\smallskip\noindent\textbf{פתרון #1.}\ }

\pagestyle{fancy}\fancyhf{}
\fancyhead[R]{\small מבוא לקוסמולוגיה --- רשימות הרצאה}
\fancyhead[L]{\small\hbox{\textdir TLT\relax\thepage}}
\renewcommand{\headrulewidth}{0.4pt}
\setlength{\parindent}{0pt}\setlength{\parskip}{4pt}

% ---- shared definitions for the homework sheets (input after ../lecture_notes/preamble) ----
\newcounter{hwq}
\newenvironment{question}[2]{%
  \refstepcounter{hwq}\needspace{5\baselineskip}%
  \begin{exbox}\textbf{\textcolor{qblue}{שאלה \thehwq\ --- #1}}\hfill\textbf{(#2 נק')}\par\smallskip}%
  {\end{exbox}}
\newenvironment{parts}{\begin{enumerate}[label=(\alph*),itemsep=3pt]}{\end{enumerate}}
\newcommand{\hwtitle}[3]{%
  \noindent\colorbox{qblue}{\parbox{\dimexpr\textwidth-2\fboxsep}{%
   \color{white}\bfseries\Large תרגיל בית #1: #2\\[3pt]\normalsize\mdseries #3}}\par\medskip}
\newcommand{\hwrules}[1]{%
\begin{goalbox}\small\textbf{הנחיות.}
\begin{itemize}
\item ארבע שאלות, 25 נקודות כל אחת. השאלות בלתי תלויות זו בזו.
\item אלא אם נאמר אחרת: $\LCDM$ שטוח עם $h=0.674$, $\Omega_m=0.315$, $\Omega_\Lambda=1-\Omega_m$ (קרינה מוזנחת), ו-$c/H_0=4448\,\Mpc$, $1/H_0=14.51\,\Gyr$.
\item רוב הסעיפים אנליטיים: גזירה, ואז הצבת מספרים במחשבון. רק סעיפים המסומנים \textbf{[מחשב]} דורשים קוד; בהם צרפו את הקוד (\en{Python} מומלץ), סמנו צירים ויחידות, ושמרו את ה-\en{seed}.
\item כל טענה, נתון או ערך שלא נלמדו בהרצאה --- עם ציטוט (מחבר, שנה, כתב עת).
#1
\end{itemize}\end{goalbox}}
\pagestyle{fancy}\fancyhf{}
\fancyhead[L]{\small\hbox{\textdir TLT\relax\thepage}}

\fancyhead[R]{\small מבוא לקוסמולוגיה --- תרגיל בית 2}
\begin{document}
\hwtitle{2}{\en{MCMC} וסופרנובות}{אחרי הרצאה 2 \quad$\cdot$\quad בייס, יסודות \en{MCMC}, דיאגרמת האבל של \en{Pantheon+}, $q_0$ ו-$H_0$}
\hwrules{\item בשאלות 2--3 מותר להשתמש בכל אלגוריתם \en{MCMC} ובכל חבילה (עם ציטוט), בתנאי שאתם מבינים מה הם עושים.
\item הנתונים לשאלות 3--4 בטבלה בסוף התרגיל.}

% ---------------------------------------------------------------------
\begin{question}{בייס עם שתי מדידות של $H_0$}{25}
\begin{parts}
\item \textbf{(7)} מדידה אחת $x$ עם שגיאה גאוסית $\sigma$ של פרמטר $\mu$, ומוקדמת גאוסית $\mu\sim\mathcal N(\mu_0,\tau^2)$. הראו שהפוסטריור גאוסי עם
\[\frac{1}{\sigma_{\rm post}^2}=\frac1{\sigma^2}+\frac1{\tau^2},\qquad \mu_{\rm post}=\sigma_{\rm post}^2\Big(\frac{x}{\sigma^2}+\frac{\mu_0}{\tau^2}\Big).\]
מה קורה בגבול $\tau\to\infty$? ובגבול $\tau\ll\sigma$?
\item \textbf{(6)} \en{Planck} נותן $H_0=67.4\pm0.5$ (דרך $\LCDM$) ו-\en{SH0ES} נותן $H_0=73.04\pm1.04\,\kmsMpc$ (סולם מרחקים). מצאו את מוטב המשולב ושגיאתו (השתמשו בסעיף (א), כשאחת המדידות משמשת "מוקדמת" לשנייה). חפשו וצטטו את שני המקורות.
\item \textbf{(6)} כמה $\sigma$ ביניהן? (רמז: ההפרש $\Delta=x_1-x_2$ גאוסי עם $\sigma_\Delta^2=\sigma_1^2+\sigma_2^2$.) מה ה-\en{p-value} הדו-צדדי? האם יש טעם לצטט את הממוצע המשולב מסעיף (ב)?
\item \textbf{(6)} נניח שב-\en{SH0ES} יש שגיאה שיטתית לא מזוהה $\sigma_{\rm sys}$ (בנוסף לשגיאה המצוטטת). כמה גדולה היא צריכה להיות כדי שהמתח ירד ל-$3\sigma$? ל-$2\sigma$? האם זה סביר בהשוואה לכיול הצפאידים ($\sim1$--$2\%$)?
\end{parts}
\end{question}

% ---------------------------------------------------------------------
\begin{question}{יסודות \en{MCMC}: מה הדוגם צריך לקיים ואיך בודקים אותו}{25}
בשאלה הזו ובשאלה 3 מותר להשתמש בכל דוגם שתבחרו (\en{Metropolis--Hastings} שכתבתם, \en{emcee}, \en{dynesty}, \en{PyMC} וכו'), עם ציטוט. הסעיפים (א)--(ג) אנליטיים.
\begin{parts}
\item \textbf{(6)} \emph{איזון מפורט.} ב-\en{Metropolis--Hastings} עם הצעה כללית $q(y|x)$ (לא בהכרח סימטרית), ההסתברות לקבל היא $\alpha=\min\!\big[1,\frac{p(y)q(x|y)}{p(x)q(y|x)}\big]$. הראו ש-$p(x)q(y|x)\alpha(x\to y)$ סימטרי ב-$x\leftrightarrow y$, ומכאן ש-$p$ היא התפלגות סטציונרית. למה מספיק לדעת את הפוסטריור עד קבוע נרמול (האבידנס)?
\item \textbf{(7)} \emph{קונטורים של גאוסיאן דו-ממדי.} פוסטריור גאוסי עם $\sigma_x=\sigma_y=1$ ומתאם $\rho=0.95$ (בדומה ל-$\Omega_m$--$\Omega_\Lambda$ בשאלה 3). (i) כתבו את $\chi^2(x,y)$ ומצאו את ההתפלגות השולית של $x$ ואת ההתפלגות המותנית $p(x|y)$. פי כמה השגיאה השולית גדולה מהמותנית? (ii) מצאו את הצירים הראשיים של האליפסה ואת אורכיהם. (iii) הראו שההסתברות בתוך $\Delta\chi^2<k$ היא $1-e^{-k/2}$ (רמז: עברו לקואורדינטות שבהן $\chi^2=r^2$), ולכן קווי $68.3\%$ ו-$95.4\%$ הם $\Delta\chi^2=2.30$ ו-$6.18$. (iv) למה ההיטל של קו ה-$68\%$ הדו-ממדי על ציר $x$ רחב מ-$\pm1\sigma$ (פי כמה)?
\item \textbf{(7)} \emph{אוטוקורלציה ומספר דגימות אפקטיבי.} מודל צעצוע לשרשרת: $x_{n+1}=\phi x_n+\sqrt{1-\phi^2}\,\sigma\xi_n$, עם $\xi_n\sim\mathcal N(0,1)$ בלתי תלויים ו-$0<\phi<1$. (i) הראו ש-$\mathrm{Var}(x_n)=\sigma^2$ במצב הסטציונרי וש-$\mathrm{Corr}(x_n,x_{n+k})=\phi^k$. (ii) הראו שעבור $N\gg1$, $\mathrm{Var}(\bar x)\simeq\sigma^2\tau/N$ עם $\tau=1+2\sum_{k\ge1}\phi^k=\frac{1+\phi}{1-\phi}$, ולכן מספר הדגימות האפקטיבי הוא $N/\tau$. (iii) עבור $\phi=0.95$: מה $\tau$, וכמה צעדים צריך כדי לדעת את הממוצע לדיוק $0.01\sigma$? מה זה אומר על "דילול" (\en{thinning}) של שרשרת?
\item \textbf{(5)} \textbf{[מחשב]} דגמו את הגאוסיאן מסעיף (ב) עם הדוגם שבחרתם. בדקו מול התוצאות האנליטיות: הסטיות השוליות, $\rho$, ואחוז הדגימות בתוך $\Delta\chi^2<2.30$ ו-$<6.18$. העריכו את $\tau$ ואת $N/\tau$ (רוב החבילות מחשבות אותם), והראו גרף \en{trace} וגרף \en{corner}.
\end{parts}
\end{question}

% ---------------------------------------------------------------------
\begin{question}{סופרנובות: היקום מאיץ \textbf{[מחשב]}}{25}
בטבלה בסוף: גדלים מתוקנים $m_B$ של 1590 סופרנובות \en{Pantheon+}, ממוצעים ב-12 סלים. המודל:
$m_{\rm th}=\mathcal M+5\log_{10}(\hat d_L/\Mpc)+25$, כאשר $\hat d_L$ הוא $d_L$ עם $H_0=70\,\kmsMpc$, ו-$\mathcal M$ היסט לא ידוע. בתרגיל הזה השתמשו בשגיאות האלכסוניות $\sigma_m$ בלבד, ומרג'נלו על $\mathcal M$ אנליטית כמו בהרצאה: $\chi^2_{\rm marg}=A-B^2/S$.
\begin{parts}
\item \textbf{(6)} חשבו את $\chi^2_{\rm marg}$ עבור ארבעה מודלים: \en{EdS} $(1,0)$, \en{Milne} $(0,0)$, פתוח $(0.3,0)$ ו-$\LCDM$ $(0.3,0.7)$ [בסוגריים $(\Omega_m,\Omega_\Lambda)$]. עם 12 נקודות, אילו מודלים נפסלים? שרטטו את השאריות $m-m_{\rm th}$ של כל המודלים יחסית ל-\en{Milne}.
\item \textbf{(6)} ב-$\LCDM$ שטוח: מצאו את $\Omega_m$ המיטבי ואת טווח $1\sigma$ ($\Delta\chi^2=1$) על גריד.
\item \textbf{(8)} הריצו \en{MCMC} (כל דוגם) על $(\Omega_m,\Omega_\Lambda)$ ללא הנחת שטיחות, עם מוקדמות אחידות $0<\Omega_m<1.5$, $-0.5<\Omega_\Lambda<1.5$ (דחו מודלים בלי מפץ גדול, שבהם $E^2(z)\le0$). הציגו \en{corner plot} עם קווי $68\%$ ו-$95\%$, דווחו חציון ו-$68\%$ ל-$\Omega_m$ ו-$\Omega_\Lambda$ ואת המתאם ביניהם. סמנו על הגרף את הקו השטוח ואת הקו $q_0=\Omega_m/2-\Omega_\Lambda=0$.
\item \textbf{(5)} מהשרשרת: מה $P(\Omega_\Lambda>0)$ ו-$P(q_0<0)$? למה הקונטור מאורך דווקא בכיוון הזה? איזו מדידה אחרת (מהרצאה 5) חותכת אותו כמעט בניצב?
\end{parts}
\end{question}

% ---------------------------------------------------------------------
\begin{question}{$q_0$ בלי מודל, ו-$H_0$ מההיסט (אנליטית)}{25}
\begin{parts}
\item \textbf{(6)} פתחו את $a(t)$ סביב $t_0$ עם $q_0=-\ddot a a/\dot a^2|_0$, והראו שעבור $z\ll1$
\[d_L=\frac{c}{H_0}\Big[z+\frac{1-q_0}{2}\,z^2+\mathcal O(z^3)\Big].\]
(מספיק יקום שטוח.) למה הפיתוח הזה לא תלוי במשוואות פרידמן או בתוכן החומרי?
\item \textbf{(7)} (i) ממשוואת התאוצה הראו ש-$q_0=\tfrac12\sum_i\Omega_i(1+3w_i)$, ולכן $q_0=\Omega_m/2-\Omega_\Lambda$; מה $q_0$ ל-\en{Planck}? (ii) בדקו את הפיתוח: פתחו את $d_L$ המדויק של \en{Milne} ($q_0=0$) ושל \en{EdS} ($q_0=\tfrac12$) עד $z^2$. (iii) בכמה מגניטודות שונה $\mu$ ב-$z=0.1$ בין $q_0=-0.53$ ל-$q_0=0$? (רמז: $\Delta\mu\simeq\frac{5}{\ln10}\frac{\Delta d_L}{d_L}$.) השוו לפיזור של סופרנובה בודדת ($\sim0.15$ מג'): כמה סופרנובות ב-$z\approx0.1$ נחוצות כדי לראות את התאוצה ב-$3\sigma$ (בהנחה ש-$\mathcal M$ ידוע)?
\item \textbf{(6)} הראו ש-$\mathcal M=M_B+5\log_{10}(70/H_0)$, כאשר $M_B$ הגודל המוחלט של סופרנובה. למה סופרנובות לבדן לא מודדות את $H_0$? הראו ש-$H_0=70\times10^{(M_B-\mathcal M)/5}\,\kmsMpc$, ושגיאה $\delta M_B$ נותנת $\delta H_0/H_0=\frac{\ln10}{5}\delta M_B$.
\item \textbf{(6)} קחו $\hat{\mathcal M}=-19.36$ (בערך מה שיוצא בשאלה 3(ב)). חשבו את $H_0$ ושגיאתו עבור (i) כיול צפאידים $M_B=-19.253\pm0.027$ (\en{SH0ES}) ו-(ii) "סולם הפוך" מ-\en{CMB+BAO}, $M_B=-19.401\pm0.027$. מצאו וצטטו את שני המקורות. מה הפרש ה-$M_B$ במגניטודות, ולאיזה הפרש במרחק (באחוזים) הוא מתאים? זהו מתח $H_0$ בשפה של סופרנובות.
\end{parts}
\end{question}

\clearpage
\begin{center}
\textbf{טבלה: סופרנובות \en{Pantheon+} ממוצעות בסלי הסחה (לשאלות 3--4).}\\[6pt]
\begin{otherlanguage}{english}
\begin{tabular}{rrrrr}
\toprule
bin & $N_{\rm SN}$ & $\bar z$ & $m_B$ [mag] & $\sigma_{m}$ [mag]\\
\midrule
1 & 154 & 0.0155 & 14.7957 & 0.0207 \\
2 & 203 & 0.0249 & 15.7940 & 0.0134 \\
3 & 150 & 0.0360 & 16.6551 & 0.0133 \\
4 & 79 & 0.0540 & 17.5313 & 0.0172 \\
5 & 44 & 0.0784 & 18.4176 & 0.0223 \\
6 & 85 & 0.1289 & 19.5458 & 0.0158 \\
7 & 261 & 0.2012 & 20.5767 & 0.0108 \\
8 & 306 & 0.3163 & 21.7074 & 0.0098 \\
9 & 179 & 0.4773 & 22.7426 & 0.0146 \\
10 & 102 & 0.6982 & 23.7163 & 0.0232 \\
11 & 19 & 1.1687 & 25.2154 & 0.0652 \\
12 & 8 & 1.7422 & 26.1644 & 0.0998 \\
\bottomrule
\end{tabular}
\end{otherlanguage}
\end{center}
{\small \textbf{הערות.} ההסחות הן $z_{\rm HD}$; השתמשו ב-$z=\bar z$ גם במודל וגם ב-$d_L$. $m_B$ ו-$\sigma_m$ הם הממוצע והשגיאה (\en{GLS}) של העקומות בכל סל, עם מטריצת השונות המשותפת הסטטיסטית${}+{}$השיטתית \emph{בתוך} הסל; מתאמים \emph{בין} סלים לא נכללו. נגזר מהשחרור הציבורי של \en{Brout et al.\ 2022, ApJ 938, 110; Scolnic et al.\ 2022, ApJ 938, 113} (עמודות \texttt{zHD}, \texttt{m\_b\_corr}, חיתוך $z_{\rm HD}>0.01$).}

\end{document}
